\section{Data and measures}\label{sec:data}

\subsection{Sample}\label{sec:sample}

The population is every captive spin-off announced in a division resolution filed with the Korean regulator between 2001 and 2026, whether or not the parent was listed. For 2017 onward we take the structured division-resolution reports that DART, the regulator's disclosure system, returns (\NNewN{} divisions, \NMixed{} of which combined a captive with an equity spin-off). For 2001 to 2016, before the structured record exists, we retrieved every filing whose title is a division or division-and-merger resolution, fetched its text, and clustered a filing and its corrections into one event. A language model then extracted the division's form and whether the new company's shares were allotted to the parent, with verbatim evidence, and we kept the \NOldN{} events with a captive part. The parents include unlisted filers, whose divisions enter the prediction tests but not the return tests. Table~\ref{tab:sample} follows the events through the three measurements we take on them.

\begin{table}[htbp]
\centering
\caption{Sample: captive spin-offs, outcome labels, prices and ownership}
\label{tab:sample}
{\footnotesize\setlength{\tabcolsep}{3pt}% generated by code/make_tables_cgir.py on 2026-10-26 — source: carveout_disposal_tracking_{2001_2016,v4}.json; dart_division_decisions_2017_2026.json; event_study_meta_compustat.csv; ownership_pre_event_all.json; gate1_pooled_20261019.json by_source
\begin{tabular}{>{\raggedright\arraybackslash}p{9.2cm}rrr}
\toprule
 & 2001--2016 & 2017--2026 & Total \\
\midrule
Captive spin-off announcements with a tracked outcome file (listed and unlisted parents) & 464 & 294 & 758 \\
\quad disposed within the tracking window & 257 & 145 & 402 \\
\quad retained at 100\% & 89 & 90 & 179 \\
\quad merged back or liquidated (excluded) & 43 & 29 & 72 \\
\quad label unresolved or too recent (excluded) & 75 & 30 & 105 \\
Clean-label events with a disposal score (Table~\ref{tab:gate1}) & 345 & 234 & 579 \\
Listed parents with daily prices around the announcement & 380 & 260 & 640 \\
\quad with a full one-year post-announcement window & 366 & 249 & 615 \\
Parents with pre-announcement controlling-block ownership & 387 & 283 & 670 \\
\bottomrule
\end{tabular}
}
\tabnote{Divisions are dated by the receipt date of the original division filing. The tracked outcome reads the parent's annual and half-year reports for up to six fiscal years after the division; ``disposed'' pools subsidiaries listed, sold by majority or diluted below 99.5 percent, and ``retained'' means 100 percent through the latest report. Clean-label events with a disposal score are those with both labels and a two-model reading of the filing. Prices: Korea Exchange adjusted prices for 2017--2026 and Compustat Global for 2001--2016; a full window has at least 90 percent of trading days $+2$ to $+250$. Ownership: controlling block at the fiscal year-end before the announcement year, from the largest-shareholder table of the annual report.}
\end{table}

The outcome is read from the parent's own later filings. For each division we traced the new company by name through the parent's annual and half-year reports from the division year for up to six fiscal years, to fiscal 2025 at the latest, and through the exchange's listing records. Two parts of the reports carry the information: the schedule of equity investments in other corporations, which the regulator returns as a table for 2017 onward, and the passages and subsidiary tables of the annual reports themselves, which a language model read with verbatim evidence for every year. We classified the subsidiary as disposed of when the parent sold a majority, listed it, or let its stake fall below 99.5 percent through a sale or a third-party investment, and as retained when the reports show 100 percent through the latest available report. Divisions whose new company was merged back into the parent or liquidated are excluded, as are divisions whose reports do not allow a classification and divisions too recent to have a record (\NAllOther{} and \NAllUnresolved{} events respectively, Table~\ref{tab:sample}). The names in the filings are provisional (``(tentative) LG Energy Solution Co., Ltd.'' in LG Chem's filing of September 2020), so each new company was also traced under the name it took at incorporation. Of the \NClean{} events with a clean label, \CleanDisposedShare{} percent were disposed of and \IPOShare{} percent were listed. \citeA{kang2025} report \KLDisposed{} percent disposed and \KLListed{} percent listed for their \KLN{} divisions of 2001--2021; the two populations were built independently. Where the label is known we also record the first fiscal year in which the parent's report shows a change in ownership, which lets us separate disposals within one year of the division from later ones.

Of the 2017--2026 labels, \NSensitive{} are sensitive to the definition or to a conflict between the two sources: a subsidiary transferred to another company of the same group, a stake temporarily diluted by an outside investor and later restored, preferred shares issued to a related party, a dilution too small to be material, and cases in which the equity-investment schedule and the annual-report text disagree. We keep the rule above for these cases and show in the Supporting Information that reclassifying all \NSensitive{} as retained, the direction least favourable to our results, leaves every result in place.

\subsection{Definitions}\label{sec:definitions}

Table~\ref{tab:defs} lists the \NDefs{} measures the paper uses, each with its operational rule and the tables that use it; the Supporting Information gives the same rows with their conceptual definitions and sources. We call the division a captive spin-off, following \citeA{kang2025}, because the two terms the US literature offers both mislead: it is not a spin-off in the sense of \citeA{miles1983}, \citeA{hite1983} and \citeA{krishnaswami1999}, since no shares are distributed, and not an equity carve-out in the sense of \citeA{schipper1986} and \citeA{nanda1991}, since no shares are sold or listed at the division. Its closest US relative is the first stage of what \citeA{klein1991} call the two stages of a carve-out, with the second stage---the sale or listing---left to the controller's discretion.

\begin{table}[htbp]
\centering
\caption{Definitions of the measures used in this paper (short form)}
\label{tab:defs}
{\footnotesize\setlength{\tabcolsep}{3pt}% generated by code/make_tables_cgir.py on 2026-10-26 — source: DEFS list in code/make_tables_cgir.py (one row list; full form S_defs_full)
\begin{tabular}{>{\raggedright\arraybackslash}p{0.6cm}>{\raggedright\arraybackslash}p{3.1cm}>{\raggedright\arraybackslash}p{7.6cm}>{\raggedright\arraybackslash}p{3.4cm}}
\toprule
No. & Term & Definition & Applied in \\
\midrule
D1 & Captive spin-off (\emph{mulchŏk bunhal}) & A division in which the parent receives all shares of the new company; the parent's shareholders receive none. & Table~\ref{tab:sample}; all tests \\ \addlinespace[1pt]
D2 & Disposal & The parent ceases to hold about 100\% of the new company within the tracking window: sale of a majority, listing, or a stake sale or third-party investment taking the stake below 99.5\%. & Tables~\ref{tab:sample}, \ref{tab:gate1}, \ref{tab:returns_by_outcome} \\ \addlinespace[1pt]
D3 & Early disposal & First ownership change recorded in the fiscal year of the division or the next. & Tables~\ref{tab:returns_by_outcome}, \ref{tab:realized_controls}, \ref{tab:block_terciles} \\ \addlinespace[1pt]
D4 & Disposal signal & The probability of disposal that two language models assign to the announcement filing, averaged. & Tables~\ref{tab:gate1}, \ref{tab:returns_by_score}, \ref{tab:controls} \\ \addlinespace[1pt]
D5 & Anonymized condition & The same filing with parent and new-company names, group prefixes and dates masked, and the filing date withheld. & Table~\ref{tab:gate1} \\ \addlinespace[1pt]
D6 & Controlling block & Common shares held by the largest shareholder and related parties, as a percentage of common shares, at the fiscal year-end before the announcement year; treasury shares excluded. & Tables~\ref{tab:block_terciles}, \ref{tab:block_regressions} \\ \addlinespace[1pt]
D7 & Announcement CAR & Market-model cumulative abnormal return over trading days $-1$ to $+1$ around the filing date. & Tables~\ref{tab:returns_by_score}, \ref{tab:controls}, \ref{tab:reform} \\ \addlinespace[1pt]
D8 & One-year BHAR & Buy-and-hold return of the parent minus that of its market index over trading days $+2$ to $+250$. & Tables~\ref{tab:returns_by_score}--\ref{tab:block_regressions} \\ \addlinespace[1pt]
D9 & Within-year rank & Percentile rank of the one-year BHAR among the year's events (0--1). & Tables~\ref{tab:controls}, \ref{tab:realized_controls}, \ref{tab:block_regressions} \\ \addlinespace[1pt]
D10 & Post-reform period & Filings received on or after 1 October 2022. & Table~\ref{tab:reform} \\ \addlinespace[1pt]
\bottomrule\end{tabular}
}
\tabnote{The same \NDefs{} rows, numbered and ordered identically, appear in full in the Supporting Information with the conceptual definition, the operational rule and the source of each.}
\end{table}

\subsection{The disposal signal read from the filing}\label{sec:signal}

We measure what a careful reader could infer from the announcement filing with two large language models, each reading the filing's text as of its date. Text-based measures of disclosure content have moved from word lists \cite{loughran2011} to supervised learning on labelled sentences \cite{frankel2022} and to language models that read words in context \cite{siano2025}. Ours differs from all three in asking the model not for a tone score but for a forecast of a specific event that we later observe, so that the measure can be validated against the record rather than against returns. Model~A is Anthropic's Claude (version \texttt{claude-fable-5-1}) and Model~B is OpenAI's GPT (version \texttt{gpt-6.1-sol}). Both were called through a cloud inference service in October 2026 with extended reasoning enabled at its highest setting and no sampling parameters altered; the system prompt, the user message template and the parsing code are in the replication archive. Each model receives the filing's full text (up to 60,000 characters) and its board-resolution date. It is instructed to forecast, from the document alone and without using anything it may know about the company's later history, the probability that the parent will not retain 100 percent of the new company within roughly three to five fiscal years. It returns that probability as an integer from 0 to 100 together with the most likely route (retention, listing, sale to a third party, a partial stake sale or investor, or an intragroup transfer or merger), three to six reasons, verbatim sentences from the filing that support them, and a set of facts stated in the filing (the stated purposes by category; whether a listing, an outside investor, a sale or buyer, or shareholder-protection measures are mentioned; the new company's business and capital). The disposal signal is the two models' probabilities averaged and rescaled to the unit interval; for 2017--2026 each model's probability is itself the mean of two independent runs. We use two models rather than one because they are trained differently and their errors need not coincide. The two models' scores correlate at \BetweenModelCorr{} across events, and each model's two runs correlate at \RepeatCorrA{} and \RepeatCorrB{} over the \NRepeat{} events judged twice, so the measure is stable.

The objection to any such measure is that the model may be recalling rather than reading: a well-known company's later listing may be in its training data. We address this in three ways. We ask each model whether it recognised the company (Model~A says it did for \ClaudeOutsideKnowledgeOld{} percent of the 2001--2016 filings, Model~B for \GPTOutsideKnowledgeOld{} percent). We re-run every filing in an anonymized condition, in the spirit of \citeA{glasserman2024}, who test the same concern for return predictions. Every occurrence of the parent's and the new company's names (with and without the corporate suffix), the names of the large business groups, and every year and date in the text are replaced by placeholders, and the board date is withheld from the message, so that recall of a named company's history is not available. And we compare the two conditions event by event. Section~\ref{sec:gate1} reports the result; if the signal survives anonymization, it is a reading of the document.

\subsection{Returns, ownership and controls}\label{sec:returns}

For listed parents we measure the announcement reaction and the return over the following year. Day 0 is the first trading day on or after the receipt date of the original division filing. The announcement return is the market-model cumulative abnormal return over days $-1$ to $+1$, with the model estimated over days $-250$ to $-31$ against the index of the parent's market (KOSPI or KOSDAQ). The one-year return is the buy-and-hold return of the parent minus that of its market index over days $+2$ to $+250$, so that the announcement reaction is excluded from it; windows with fewer than 90 percent of their trading days are dropped, so a parent delisted early in the year leaves the one-year sample. Daily adjusted prices for 2017--2026 come from the Korea Exchange, with halt days removed. For 2001--2016, for which the exchange's adjusted series are not available, they come from Compustat Global's daily security file, which carries delisted companies; we verified its returns against the exchange's on the 2017--2026 overlap (median daily-return correlation \KRXCompustatCorr{} across \NPriceOverlap{} events). Because Compustat occasionally misses a corporate-action adjustment, we exclude from 2001--2016 the parents whose window contains a daily return beyond the exchange's price limit without a delisting (\GlitchExcluded{} events; the Supporting Information shows the results with them included). Long-run returns are skewed, and buy-and-hold tests are known to over-reject when events cluster in time \cite{barber1997, lyon1999, fama1998}. Beside means we therefore report medians, winsorised means, bootstrap confidence intervals and the parent's percentile rank within its announcement year, and we add two calendar-time tests of the kind \citeA{mitchell2000} recommend: an equal-weighted daily portfolio of all parents within their one-year windows, and a monthly portfolio regressed on Korean five-factor returns.

The controlling block is the percentage of the parent's common shares held by the largest shareholder and the related parties the filing lists with it (family members, affiliated companies, officers), excluding treasury shares, at the fiscal year-end before the announcement year. For announcements from 2016 onward we take it from the largest-shareholder table of the annual report through the regulator's interface. For earlier announcements, which that interface does not cover, we read the same table from the annual report itself, using the pre-division year's report where we have it and otherwise the year-start column of the division year's report. Two models extracted the figures independently from the table text and agree within one percentage point in \OwnWithinOne{} of \OwnPairs{} cases; we drop the \OwnDisagree{} disagreements (Supporting Information). The median block is \BlockMedian{} percent and the median largest shareholder alone holds \LargestMedian{} percent; \citeA{kang2025}, whose measure is the controlling family's and inside directors' stake, report a mean of \KLOwnMean{} and a median of \KLOwnMedian{} percent.

The controls are the parent's market capitalisation at the month-end before the announcement and its book-to-market ratio (Compustat Global), its market-adjusted return over the six months before the announcement, and its market. Three further controls are stated in the filing: the new company's assets relative to the surviving company's, the surviving company's debt-to-assets ratio, and whether the filing states a listing application. For 2017--2026 we add the parent's return on assets, debt ratio and size in the fiscal year before the division from the regulator's financial-statement interface, which does not reach earlier years. Missing controls are set to the sample median with an indicator, so that the sample does not shrink with control availability.
